\begin{mistake}{2026-04-12}{28}

\begin{problemcontent}
设 \( y=y(x) \) 是一阶线性方程 \( \dfrac{\mathrm{d}y}{\mathrm{d}x}+P(x)y=Q(x) \) 的解，其中 \( P(x),Q(x) \) 都以 \( T \) 为周期，且在 \(( -\infty,+\infty )\) 上连续。证明：\( y(0)=y(T) \) 是 \( y(x) \) 以 \( T \) 为周期的充要条件。
\end{problemcontent}

\begin{solutioncontent}
证明分为必要性与充分性两部分。

必要性：若 \(y(x)\) 以 \(T\) 为周期，则对任意 \(x\) 都有
\[
y(x+T)=y(x)
\]
令 \(x=0\)，立即得到
\[
y(T)=y(0)
\]
故条件是必要的。

充分性：现设 \(y(0)=y(T)\)。构造函数
\[
u(x)=y(x+T)
\]
则
\[
u'(x)=y'(x+T)
\]
由于 \(y\) 满足原方程，故
\[
y'(x+T)+P(x+T)y(x+T)=Q(x+T)
\]
又因为 \(P,Q\) 都以 \(T\) 为周期，所以
\[
P(x+T)=P(x),\qquad Q(x+T)=Q(x)
\]
从而
\[
u'(x)+P(x)u(x)=Q(x)
\]
这表明 \(u(x)\) 与 \(y(x)\) 满足同一个一阶线性微分方程。

再看初值：
\[
u(0)=y(T)=y(0)
\]
因此 \(u\) 与 \(y\) 满足同一微分方程且具有相同初值。由于 \(P(x),Q(x)\) 在全轴上连续，一阶线性初值问题解唯一，所以
\[
u(x)\equiv y(x)
\]
即
\[
y(x+T)=y(x)
\]
故 \(y(x)\) 以 \(T\) 为周期。

综上，\(y(0)=y(T)\) 是 \(y(x)\) 以 \(T\) 为周期的充要条件。
\end{solutioncontent}

\mistakeknowledge{
\begin{itemize}
  \item \textbf{核心公式}：一阶线性初值问题在系数连续时解存在且唯一；证明周期性常用平移构造 \(u(x)=y(x+T)\)。
  \item \textbf{易错点}：仅由 \(y(0)=y(T)\) 不能直接推出周期性，关键步骤是证明平移后的函数满足同一方程与同一初值，再用唯一性定理。
  \item \textbf{关联知识}：对周期系数微分方程，判断某个解是否为周期解，常转化为“平移后是否仍是同一初值问题”的判定。
\end{itemize}}

\end{mistake}
